\documentclass{article}

\usepackage[preprint]{neurips_2026}

\usepackage[utf8]{inputenc}
\usepackage[T1]{fontenc}
\usepackage{hyperref}
\usepackage{url}
\usepackage{booktabs}
\usepackage{amsfonts}
\usepackage{nicefrac}
\usepackage{microtype}
\usepackage{xcolor}
\usepackage{enumitem}

\title{KISS Sorcar: A Relentless, Self-Continuing AI Agent \\
  for Integrated Software Development}

\author{%
  Koushik Sen \\
  Department of Electrical Engineering and Computer Sciences \\
  University of California, Berkeley \\
  Berkeley, CA 94720 \\
  \texttt{ksen@berkeley.edu} \\
}

\begin{document}

\maketitle

\begin{abstract}
We present KISS Sorcar, an AI-powered general assistant and integrated
development environment built around a layered agent architecture.  At
its foundation, the KISS Agent implements a budget-aware ReAct loop
with native function calling.  The Relentless Agent extends this with
automatic multi-session continuation and stall detection, enabling
arbitrarily long tasks that survive context-window exhaustion.  The
Sorcar Agent adds coding tools, browser automation, and parallel
sub-agent execution.  The Stateful Sorcar Agent introduces persistent
multi-turn chat sessions, and the Worktree Sorcar Agent isolates every
task on a dedicated git branch via worktrees so the developer's main
working tree is never modified.  We describe the design of each layer,
the structured system prompt that governs agent behavior, and the
engineering decisions that make the system practical for real-world
software development.
\end{abstract}

%----------------------------------------------------------------------
\section{Introduction}
\label{sec:intro}
%----------------------------------------------------------------------

Large language models (LLMs) have demonstrated remarkable ability to
generate code, reason about software architecture, and operate
developer tools~\citep{chen2021evaluating,roziere2023code}.  However,
using a raw LLM as a software engineering assistant exposes several
practical gaps: context windows are finite, a single mistake can derail
an entire session, and generated changes are difficult to review or
revert once applied to a live codebase.

KISS Sorcar addresses these gaps through a five-layer agent hierarchy
in which each layer adds exactly one concern:

\begin{enumerate}[nosep]
  \item \textbf{KISS Agent} --- budget-tracked ReAct loop with native
    function calling.
  \item \textbf{Relentless Agent} --- automatic continuation across
    sub-sessions with stall detection.
  \item \textbf{Sorcar Agent} --- coding tools, browser automation,
    and parallel sub-agent execution.
  \item \textbf{Stateful Sorcar Agent} --- persistent multi-turn chat
    sessions with history recall.
  \item \textbf{Worktree Sorcar Agent} --- git-worktree isolation so
    every task runs on its own branch.
\end{enumerate}

The name ``KISS'' reflects the Keep It Simple, Stupid design
philosophy: each layer is small, each concern is isolated, and the
overall system avoids unnecessary abstraction.  The name ``Sorcar''
pays homage to P.~C.~Sorcar, the legendary Indian magician, evoking
the idea of an agent that performs feats that appear magical yet are
grounded in disciplined engineering.

The system is complemented by a carefully structured system prompt that
encodes invariants about code quality, verification discipline, and
self-improvement.  Together, the architecture and the prompt form a
practical AI IDE that integrates with Visual Studio Code and operates
on real-world repositories.

%----------------------------------------------------------------------
\section{Motivation and Key Design Principles}
\label{sec:motivation}
%----------------------------------------------------------------------

\paragraph{Finite context, infinite tasks.}
Even the largest context windows are insufficient for complex
refactoring or multi-file bug hunts.  Rather than hoping the model
finishes before hitting the limit, the system is designed from the
ground up to \emph{continue} across multiple sub-sessions, each with
its own context window, while preserving a compressed summary of prior
progress.

\paragraph{Budget accountability.}
Every API call costs money.  The agent tracks input tokens, output
tokens, cache hits, and cache writes at every step, computes a running
dollar cost, and enforces per-task and global budget ceilings.  If the
budget is exceeded, the agent raises a hard error rather than silently
accumulating charges.

\paragraph{Safe-by-default code changes.}
Applying AI-generated changes directly to a developer's working tree
is risky: a bad edit can break a build, and reverting requires manual
git archaeology.  The worktree layer eliminates this risk entirely by
running every task on a disposable git branch in a separate worktree
directory.  The developer reviews changes and explicitly chooses to
merge or discard.

\paragraph{No mocks, no fakes, real integration tests.}
The system's own test suite enforces a strict policy: tests must
exercise actual behavior, not mocked approximations.  This philosophy
extends to the agent's instructions --- when the agent writes tests
for user code, it is instructed to avoid mocks and achieve full branch
coverage.

\paragraph{Relentless execution.}
The system is designed to be \emph{relentless}: it does not give up
when a sub-session fails.  Instead, it summarizes what happened,
detects whether the same failure is recurring, and either continues
with a fresh strategy or escalates to the user.

%----------------------------------------------------------------------
\section{Agent Architecture}
\label{sec:architecture}
%----------------------------------------------------------------------

The five agent layers form a strict inheritance chain.  Each layer
delegates upward for the concerns it does not own.

%......................................................................
\subsection{KISS Agent}
\label{sec:kiss-agent}
%......................................................................

The KISS Agent is the innermost execution unit.  It implements a
standard ReAct loop~\citep{yao2023react} with the following
characteristics:

\textbf{Native function calling.}  Tools are registered as ordinary
Python callables.  The agent builds an OpenAI-compatible tool schema
once at setup time and caches it, avoiding redundant schema
construction on every LLM call.  A special \texttt{finish} tool
signals task completion and returns the result to the caller.

\textbf{Step, token, and budget tracking.}  At every step the agent
extracts input and output token counts from the API response, computes
the dollar cost using a per-model pricing table, and updates both a
local budget counter and a global (cross-agent) budget counter
protected by a class-level lock.  The agent checks three limits before
each step: per-agent budget, global budget, and maximum step count.

\textbf{Error resilience.}  Transient API errors (rate limits, server
errors) are retried up to a configurable threshold of consecutive
failures.  Non-retryable errors (authentication failures, permission
denials) are detected by matching error type names and message
substrings, and are raised immediately.

\textbf{Non-agentic mode.}  When tools are not needed, the agent can
run a single generation without the ReAct loop, which is useful for
summarization or question-answering sub-tasks.

The KISS Agent is stateless across runs: each call to its run method
resets the conversation, token counters, and tool registry.  This
makes it safe to reuse a single agent instance for multiple sequential
tasks.

%......................................................................
\subsection{Relentless Agent}
\label{sec:relentless-agent}
%......................................................................

The Relentless Agent wraps a KISS Agent in a continuation loop.  Its
core contribution is the ability to execute tasks that exceed a single
context window by breaking them into sub-sessions.

\textbf{Continuation protocol.}  The \texttt{finish} tool exposed to
the inner KISS Agent accepts three fields: a success flag, a continue
flag, and a summary.  When the agent sets \texttt{is\_continue=True},
the Relentless Agent starts a new sub-session with a fresh context
window.  The prompt for the new session includes a chronologically
ordered list of all prior attempt summaries, instructing the agent not
to redo completed work.

\textbf{Forced continuation on failure.}  If a sub-session raises an
exception (e.g., the step limit is hit before the agent calls finish),
the Relentless Agent does not abort.  Instead, it saves the full
trajectory to a temporary file, spawns a separate summarizer agent to
read the trajectory and produce a concise summary, and then uses that
summary as the progress text for the next sub-session.  This ensures
that even crashed sessions contribute useful context to subsequent
attempts.

\textbf{Stall detection.}  After each continuation, the agent extracts
normalized error phrases from the summary and compares them across the
most recent sub-sessions.  If the same error phrases appear in three
consecutive summaries, the system declares a stall and terminates with
a diagnostic message rather than wasting budget on an infinite retry
loop.  A stall warning is injected into the prompt after three
attempts, instructing the agent to abandon its current strategy and
try a fundamentally different approach.

\textbf{Budget aggregation.}  Token counts and dollar costs are
accumulated across all sub-sessions.  Each new sub-session receives
offsets so that the usage display shown to the user reflects the true
cumulative cost, not just the cost of the current sub-session.

%......................................................................
\subsection{Sorcar Agent}
\label{sec:sorcar-agent}
%......................................................................

The Sorcar Agent adds the tools that make the system useful for
software development and general-purpose automation.

\textbf{Coding tools.}  Four core tools are provided: a shell command
executor with streaming output, a file reader, a precise string-based
file editor, and a file writer.  The shell executor supports a
configurable timeout, streams output to the user interface in real
time, and respects a stop event that allows the user to cancel a
running command.

\textbf{Browser automation.}  A web-use tool provides programmatic
browser control: navigating to URLs, reading page accessibility trees,
clicking elements, typing text, pressing keys, scrolling, and taking
screenshots.  This enables the agent to research documentation, verify
deployed applications, and interact with web-based tools.

\textbf{Parallel sub-agents.}  An optional parallel execution tool
spawns independent Sorcar Agent instances in a thread pool.  Each
sub-agent gets its own LLM context and tool set.  This is useful for
embarrassingly parallel tasks such as summarizing multiple files or
researching independent topics.  The results are collected and
returned in input order.

\textbf{User interaction.}  An ask-user-question tool allows the agent
to pause execution and request clarification from the user.  In the
VS~Code integration, this renders as a text input in the sidebar; in
CLI mode, it reads from standard input.

\textbf{Docker isolation.}  When a Docker image is specified, the
coding tools are replaced with Docker-aware variants that execute
commands inside a container, providing an additional layer of
sandboxing for untrusted tasks.

%......................................................................
\subsection{Stateful Sorcar Agent}
\label{sec:stateful-agent}
%......................................................................

The Stateful Sorcar Agent adds multi-turn conversation persistence.

\textbf{Chat sessions.}  Each task is assigned to a chat session
identified by a stable chat ID.  The agent persists tasks and their
results to a local database.  When a new task arrives within the same
chat session, the agent loads all prior tasks and results and prepends
them to the prompt as numbered context entries, allowing the LLM to
reference earlier work.

\textbf{Session management.}  The agent supports three operations:
starting a new chat (which clears the chat ID), resuming a chat by
task description (which looks up the corresponding chat ID), and
resuming by explicit chat ID.  This enables both automatic session
continuity in an IDE and manual session selection from the command
line.

\textbf{Metadata persistence.}  After each task, the agent records
metadata including the model used, working directory, software version,
token count, cost, and whether the task used parallel execution or
worktree isolation.  This audit trail supports cost analysis and
debugging.

%......................................................................
\subsection{Worktree Sorcar Agent}
\label{sec:worktree-agent}
%......................................................................

The Worktree Sorcar Agent is the outermost layer and the one that
users interact with in the VS~Code extension.  Its defining feature is
git-worktree isolation.

\textbf{Branch-per-task.}  When a task starts, the agent creates a new
git branch and a corresponding worktree directory.  The branch name
encodes the chat ID and a timestamp for uniqueness.  All agent
modifications happen inside the worktree; the user's main working tree
remains untouched.

\textbf{Dirty-state preservation.}  If the user's main working tree
has uncommitted changes, the agent copies them into the worktree and
creates a baseline commit.  This ensures the agent sees the same state
the user sees, while keeping the user's actual index and working tree
clean.  During merge, the system uses cherry-pick from the baseline
commit to replay only the agent's changes, excluding the
dirty-state snapshot.

\textbf{Auto-commit and auto-merge.}  When the user starts a new task
or a new chat, any pending worktree from a previous task is
automatically committed with an LLM-generated commit message and
squash-merged into the original branch.  If the merge has conflicts,
the branch is preserved and the user receives instructions for manual
resolution.

\textbf{Concurrency safety.}  A per-repository file lock serializes
the checkout, stash, merge, and pop sequence so that concurrent tabs
in the IDE cannot interleave operations on the same repository.
Thread-local storage isolates per-task state (stream parsing buffers,
bash output buffers, recording state) so that stopping one task does
not corrupt another.

\textbf{Crash recovery.}  All worktree state is stored in git itself
(branch names, git config entries) rather than in sidecar files.  On
process restart, the agent queries git for any pending
branch matching its chat ID prefix and reconstructs all instance
attributes from git config, enabling seamless recovery.

\textbf{Graceful fallback.}  If the working directory is not inside a
git repository, if the repository has no commits, or if HEAD is
detached, the agent falls back to direct execution without worktree
isolation, ensuring it never fails due to git preconditions.

%----------------------------------------------------------------------
\section{The System Prompt}
\label{sec:system-prompt}
%----------------------------------------------------------------------

The system prompt is a structured document that governs the agent's
behavior across all tasks.  It is not a generic instruction to ``be
helpful'' but rather a precise specification of engineering discipline.
We describe its key sections.

%......................................................................
\subsection{Execution Mindset}
%......................................................................

The prompt opens with a directive that sets the tone: \emph{``Focus on
the given task.  Its completion is your sole goal.  Be relentless.  Be
calm.  Be rigorous.  Be accurate.  Check facts.  No AI slop.''}  This
framing discourages the model from hedging, apologizing, or producing
filler text, and encourages decisive, verification-oriented behavior.

%......................................................................
\subsection{Identity and Tool Rules}
%......................................................................

The agent is given a concrete identity (name, developer, repository
URL) rather than a vague persona.  Tool usage rules are explicit: use
the write tool for new files and the edit tool for small changes; run
shell commands synchronously with a timeout; use the browser tool for
web navigation; call finish when done.  These rules eliminate ambiguity
about which tool to use in which situation.

A rule about the \texttt{PageDrop} API allows the agent to upload HTML
files and share them via a public URL, enabling it to create and share
rich reports, visualizations, or documentation.

%......................................................................
\subsection{Pre-flight Checks}
%......................................................................

Before modifying any file, the agent is instructed to read it first.
If a task depends on existing architecture or behavior, the agent must
read the relevant source files.  If referenced files, commands, or
configuration do not exist, the agent must stop and ask rather than
guessing.  When fixing a bug, the agent must write a test that
reproduces the bug before attempting a fix.

These pre-flight checks prevent a common failure mode in which the
agent modifies a file based on an incorrect assumption about its
contents, introducing subtle bugs that are difficult to diagnose.

%......................................................................
\subsection{Code Style Guidelines}
%......................................................................

The prompt encodes a minimalist code philosophy: write simple, clean,
readable code with minimal indirection; avoid unnecessary attributes,
variables, and configuration; avoid tight coupling; eliminate
duplication; document all public methods.  Critically, the agent is
instructed to understand the root cause of a bug and patch the root
cause, not apply a superficial fix.

A metacognitive instruction asks the agent to pause before writing
code and consider whether the planned code is simple, elegant,
general, and minimal.  This encourages the model to spend more
inference-time compute on design before committing to an
implementation.

%......................................................................
\subsection{Deep Work Rules}
%......................................................................

When a task says ``align'', ``match'', or ``make consistent'', the
agent must first read the target to determine the exact desired state,
then write specific concrete values.  It must never edit based on a
vague reference.  For multi-part work, it must list planned changes
before executing.  Every meaningful change must have a concrete
verification method.

These rules address a failure mode where the model interprets an
instruction loosely and makes changes that are directionally correct
but concretely wrong.

%......................................................................
\subsection{Planning for Complex Tasks}
%......................................................................

For tasks involving three or more files or cross-module changes, the
agent must produce an explicit plan: list the files that need to
change and why, state the exact intended change in each file, identify
dependencies and execution order, and state how each change will be
verified.  For simple single-file tasks, formal planning is skipped.

This adaptive planning rule avoids both the overhead of planning
trivial changes and the risk of ad-hoc execution on complex changes.

%......................................................................
\subsection{Testing Instructions}
%......................................................................

The testing section is perhaps the most opinionated: the agent must
achieve 100\% branch coverage, tests must not use mocks, patches,
fakes, or any form of test doubles, all tests must be integration
tests, and each test must be independent and verify actual behavior.

The prohibition on mocks is a deliberate design choice.  Mocks test
that code calls certain methods in a certain order --- they test the
implementation, not the behavior.  Integration tests verify that the
system actually works, which is what matters for a tool that modifies
real codebases.

%......................................................................
\subsection{Self-Improvement Loop}
%......................................................................

The agent maintains a preferences file that captures user invariants
discovered during task execution.  At the start of each task, it reads
this file; at the end, it updates it with any new preferences
inferred from the task.  Conflicting preferences are resolved in favor
of the newer observation.

This mechanism allows the agent to learn persistent user preferences
(e.g., preferred code style, testing conventions, project-specific
invariants) without requiring the user to repeat them in every prompt.

%......................................................................
\subsection{Pre-Finish Verification}
%......................................................................

Before declaring a task complete, the agent must: re-read every
modified file and verify correctness, run lint, type checking, and
tests, explicitly check each user requirement against what was
delivered, and continue working if any check fails.  If the same fix
has been retried three times without progress, the agent must step
back and try a different strategy.

This verification protocol is the system-prompt analog of the
Relentless Agent's stall detection: it prevents the agent from
declaring premature success and forces it to confront failures.

%----------------------------------------------------------------------
\section{Related Work}
\label{sec:related}
%----------------------------------------------------------------------

\textbf{Code generation agents.}
SWE-Agent~\citep{yang2024sweagent} and
OpenHands~\citep{wang2024openhands} provide LLM-based agents for
software engineering tasks.  Both use a single-session execution model
without automatic continuation.  KISS Sorcar's Relentless Agent layer
addresses the single-session limitation directly.

\textbf{ReAct and tool use.}
The ReAct framework~\citep{yao2023react} interleaves reasoning and
action.  Toolformer~\citep{schick2023toolformer} teaches models to use
tools via self-supervised learning.  KISS Sorcar uses native function
calling provided by modern LLM APIs rather than in-context tool
descriptions, which reduces prompt overhead and improves reliability.

\textbf{Multi-turn agents.}
AutoGPT~\citep{autogpt2023} and BabyAGI~\citep{babyagi2023} implement
multi-step autonomous agents.  These systems typically lack budget
controls, stall detection, and safe code isolation.  KISS Sorcar's
layered architecture addresses each of these concerns in a dedicated
layer.

\textbf{IDE integration.}
GitHub Copilot~\citep{copilot2021}, Cursor~\citep{cursor2024}, and
Windsurf~\citep{windsurf2024} provide AI assistance within editors.
These tools typically operate at the level of code completion or
single-turn chat.  KISS Sorcar operates at the level of autonomous
multi-step task execution with full tool access and git-level
isolation.

%----------------------------------------------------------------------
\section{Conclusion}
\label{sec:conclusion}
%----------------------------------------------------------------------

KISS Sorcar demonstrates that a layered, single-concern architecture
can address the practical challenges of deploying LLM agents for
real-world software development.  The KISS Agent provides a robust
execution core with budget tracking.  The Relentless Agent eliminates
the context-window ceiling through automatic continuation and stall
detection.  The Sorcar Agent equips the agent with the tools a
developer actually uses.  The Stateful Sorcar Agent enables multi-turn
conversations that persist across sessions.  The Worktree Sorcar Agent
ensures that every task is safely isolated in git, giving the
developer full control over what gets merged.

The system prompt complements the architecture by encoding engineering
discipline directly into the agent's behavior: read before writing,
test before fixing, plan before executing, verify before finishing.
These are not novel insights --- they are the practices of careful
software engineering, translated into instructions that an LLM can
follow.

The source code is publicly available at
\url{https://github.com/ksenxx/kiss_ai}.

%----------------------------------------------------------------------
% References
%----------------------------------------------------------------------

{
\small

\bibliographystyle{plainnat}

\begin{thebibliography}{99}

\bibitem[Chen et~al.(2021)]{chen2021evaluating}
Chen, M., Tworek, J., Jun, H., Yuan, Q., Pinto, H.~P., Kaplan, J.,
Edwards, H., Burda, Y., Joseph, N., Brockman, G., et~al.
\newblock Evaluating large language models trained on code.
\newblock {\em arXiv preprint arXiv:2107.03374}, 2021.

\bibitem[Rozi\`{e}re et~al.(2023)]{roziere2023code}
Rozi\`{e}re, B., Gehring, J., Gloeckle, F., Sootla, S., Gat, I.,
Tan, X.~E., Adi, Y., Liu, J., Remez, T., Rapin, J., et~al.
\newblock Code Llama: Open foundation models for code.
\newblock {\em arXiv preprint arXiv:2308.12950}, 2023.

\bibitem[Yao et~al.(2023)]{yao2023react}
Yao, S., Zhao, J., Yu, D., Du, N., Shafran, I., Narasimhan, K., and
Cao, Y.
\newblock ReAct: Synergizing reasoning and acting in language models.
\newblock In {\em International Conference on Learning
  Representations (ICLR)}, 2023.

\bibitem[Schick et~al.(2023)]{schick2023toolformer}
Schick, T., Dwivedi-Yu, J., Dess\`{i}, R., Raileanu, R., Lomeli, M.,
Zettlemoyer, L., Cancedda, N., and Scialom, T.
\newblock Toolformer: Language models can teach themselves to use
  tools.
\newblock In {\em Advances in Neural Information Processing Systems
  (NeurIPS)}, 2023.

\bibitem[Yang et~al.(2024)]{yang2024sweagent}
Yang, J., Jimenez, C.~E., Wettig, A., Liber, K., Narasimhan, K., and
Press, O.
\newblock {SWE-agent}: Agent-computer interfaces enable automated
  software engineering.
\newblock {\em arXiv preprint arXiv:2405.15793}, 2024.

\bibitem[Wang et~al.(2024)]{wang2024openhands}
Wang, X., Chen, Y., Yuan, L., Zhang, Y., Li, Y., Peng, H., and Ji,
H.
\newblock {OpenHands}: An open platform for {AI} software developers as
  generalist agents.
\newblock {\em arXiv preprint arXiv:2407.16741}, 2024.

\bibitem[{Significant Gravitas}(2023)]{autogpt2023}
{Significant Gravitas}.
\newblock {AutoGPT}.
\newblock \url{https://github.com/Significant-Gravitas/AutoGPT}, 2023.

\bibitem[Nakajima(2023)]{babyagi2023}
Nakajima, Y.
\newblock {BabyAGI}.
\newblock \url{https://github.com/yoheinakajima/babyagi}, 2023.

\bibitem[{GitHub}(2021)]{copilot2021}
{GitHub}.
\newblock {GitHub Copilot}: Your {AI} pair programmer.
\newblock \url{https://github.com/features/copilot}, 2021.

\bibitem[{Cursor}(2024)]{cursor2024}
{Cursor}.
\newblock {Cursor}: The {AI}-first code editor.
\newblock \url{https://cursor.sh}, 2024.

\bibitem[{Codeium}(2024)]{windsurf2024}
{Codeium}.
\newblock {Windsurf}: The {AI}-powered {IDE}.
\newblock \url{https://windsurf.com}, 2024.

\end{thebibliography}
}

%%%%%%%%%%%%%%%%%%%%%%%%%%%%%%%%%%%%%%%%%%%%%%%%%%%%%%%%%%%%
\newpage
\input{checklist.tex}

\end{document}
